% 付録：行列の基本

\section{行列の基本}\label{節：行列の基本}

本編では,複雑な\bookterm{recurrence}{漸化式}を等比型などの既知の形へ変えることを繰り返してきた.
付録では,その置き換えがなぜうまくいくのかを,\bookterm{matrix}{行列}や\bookterm{difference-operator}{差分},\bookterm{generating}{母関数}という言葉で捉え直す.
最初に,そのための道具となる\bookterm{matrix}{行列}の読み方と計算の仕方を説明しよう.
\bookterm{matrix}{行列}を初めて学ぶ読者も,数値例から順に確かめてほしい.
ここでは主に二行二列の\bookterm{matrix}{行列}を扱う.
後の節では,本編で求めた\bookterm{fixed}{不動点}や\bookterm{geometric}{等比数列}が,ここで学ぶ計算の中にそのまま現れる.

\subsection{行列とは何か}
  \BookFigMatrixMap


\bookterm{matrix}{行列}とは,数を長方形に並べ,全体を括弧で囲んだものである.
たとえば
\[
  A=\begin{pmatrix}2&1\\1&2\end{pmatrix}
\]
は,縦に二段,横に二列の数を並べた\bookterm{matrix}{行列}である.
横の並びを\textbf{行},縦の並びを\textbf{列}という.
上が第一行,下が第二行であり,左が第一列,右が第二列である.
並んでいる一つ一つの数を\textbf{成分}と呼ぶ.
この例の第一行・第二列の成分は$1$である.

また,二つの数を縦に並べた
\[
  \boldsymbol{v}=\begin{pmatrix}x\\y\end{pmatrix}
\]
を\textbf{\bookterm{column-vector}{列ベクトル}}と呼ぶ.
太字の$\boldsymbol{v}$は,二つの数をまとめて表すための記号である.
平面上の点$(x,y)$への矢印と考えることもできるが,以下ではまず「二つの値を入れた一組の記録」として読めばよい.

\bookterm{column-vector}{列ベクトル}の足し算と定数倍は,成分ごとに行う.
\[
  \begin{pmatrix}x\\y\end{pmatrix}
  +\begin{pmatrix}u\\v\end{pmatrix}
  =\begin{pmatrix}x+u\\y+v\end{pmatrix},
  \qquad
  c\begin{pmatrix}x\\y\end{pmatrix}
  =\begin{pmatrix}cx\\cy\end{pmatrix}.
\]
たとえば$\begin{pmatrix}3\\1\end{pmatrix}+\begin{pmatrix}1\\-1\end{pmatrix}
=\begin{pmatrix}4\\0\end{pmatrix}$である.
二つの\bookterm{column-vector}{列ベクトル}が等しいとは,上下それぞれの成分が等しいことをいう.
一般の二成分のベクトルも
\[
  \begin{pmatrix}x\\y\end{pmatrix}
  =\frac{x+y}{2}\begin{pmatrix}1\\1\end{pmatrix}
   +\frac{x-y}{2}\begin{pmatrix}1\\-1\end{pmatrix}
\]
と分けられる.上下の成分をそれぞれ足して確かめてほしい.
本編の\sectionref*{節：連立型}で和と差を求めた操作は,この二つの係数を求めることに対応する.
この対応は\bookproblemref{practice-basic}{7}{標準問7}で確認できる.
後で,この二つの方向が更新の定数倍とどう結び付くかを見る.

\subsubsection{行列を列ベクトルに掛ける}

先ほどの\bookterm{matrix}{行列}$A$を\bookterm{column-vector}{列ベクトル}に掛けるときは,次のように計算する.
\[
  \begin{pmatrix}2&1\\1&2\end{pmatrix}
  \begin{pmatrix}x\\y\end{pmatrix}
  =
  \begin{pmatrix}2x+y\\x+2y\end{pmatrix}.
\]
上の成分は,第一行の$2,\,1$を$x,y$にそれぞれ掛けて足したもの.
下の成分は,第二行の$1,\,2$を$x,y$にそれぞれ掛けて足したものである.
たとえば
\[
  A\begin{pmatrix}1\\0\end{pmatrix}
  =\begin{pmatrix}2\\1\end{pmatrix},
  \qquad
  A\begin{pmatrix}2\\1\end{pmatrix}
  =\begin{pmatrix}5\\4\end{pmatrix}.
\]
この\bookterm{matrix}{行列}は,二つの値$x,y$を受け取り,新しい二つの値$2x+y,x+2y$を作る規則を表している.

一般には
\[
  \begin{pmatrix}p&q\\r&s\end{pmatrix}
  \begin{pmatrix}x\\y\end{pmatrix}
  =\begin{pmatrix}px+qy\\rx+sy\end{pmatrix}
\]
と定める.
したがって,連立した二つの式
\[
  x'=px+qy,\qquad y'=rx+sy
\]
を,\bookterm{matrix}{行列}と\bookterm{column-vector}{列ベクトル}の積一つにまとめられる.
ここで$x',y'$は,更新後の値を表す記号である.

また,成分ごとに展開すれば
\[
  A(c\boldsymbol{u}+d\boldsymbol{v})
  =cA\boldsymbol{u}+dA\boldsymbol{v}
\]
が成り立つ.
これは,二つの\bookterm{column-vector}{列ベクトル}を足してから更新しても,別々に更新してから足しても結果が同じ,という性質である.
後で初期値を二つの部分に分けて考える際に使う.

\subsection{行列の積は,更新を続けて行うことを表す}

\bookterm{matrix}{行列}をもう一度掛けると,同じ規則で二回更新できる.
先ほどの例なら
\[
  \begin{pmatrix}1\\0\end{pmatrix}
  \ \longmapsto\
  \begin{pmatrix}2\\1\end{pmatrix}
  \ \longmapsto\
  \begin{pmatrix}5\\4\end{pmatrix}
\]
と進む.
この二回分の更新を一つにまとめた\bookterm{matrix}{行列}を$A^2$と書く.
実際に文字$x,y$で計算すると
\[
  \begin{aligned}
    A\left(A\begin{pmatrix}x\\y\end{pmatrix}\right)
    &=A\begin{pmatrix}2x+y\\x+2y\end{pmatrix}\\
    &=\begin{pmatrix}5x+4y\\4x+5y\end{pmatrix}.
  \end{aligned}
\]
したがって
\[
  A^2=\begin{pmatrix}5&4\\4&5\end{pmatrix}
\]
である.
各成分を単に二乗するのではなく,更新を二回行った結果に合わせて積を定めている.

異なる\bookterm{matrix}{行列}$A,B$でも,まず$B$,次に$A$を掛ける操作を$AB$と書く.
$B$の各列に$A$を掛けた結果を並べると,積$AB$が得られる.
具体的には
\[
  \begin{aligned}
    &\begin{pmatrix}p&q\\r&s\end{pmatrix}
     \begin{pmatrix}u&v\\w&z\end{pmatrix}\\
    &\qquad=
     \begin{pmatrix}
       pu+qw&pv+qz\\
       ru+sw&rv+sz
     \end{pmatrix}.
  \end{aligned}
\]
右の\bookterm{matrix}{行列}から先に作用するため,一般には$AB$と$BA$は一致しない.

何も変えない更新は
\[
  I=\begin{pmatrix}1&0\\0&1\end{pmatrix},
  \qquad
  I\begin{pmatrix}x\\y\end{pmatrix}
  =\begin{pmatrix}x\\y\end{pmatrix}
\]
で表される.
$I$を\textbf{\bookterm{identity}{単位行列}}という.
数の累乗で$a^0=1$とするのに対応して,\bookterm{matrix}{行列}では$A^0=I$と定める.

\subsection{固有ベクトルと固有値：更新が定数倍になる場合}
  \BookFigEigen


毎回\bookterm{matrix}{行列}の積を計算するのは手間がかかる.
そこで,\bookterm{matrix}{行列}を掛けても定数倍になるだけの\bookterm{column-vector}{列ベクトル}を探してみよう.
先ほどの$A$について
\[
  \begin{aligned}
    A\begin{pmatrix}1\\1\end{pmatrix}
      &=\begin{pmatrix}3\\3\end{pmatrix}
       =3\begin{pmatrix}1\\1\end{pmatrix},\\
    A\begin{pmatrix}1\\-1\end{pmatrix}
      &=\begin{pmatrix}1\\-1\end{pmatrix}
  \end{aligned}
\]
である.
第一の\bookterm{column-vector}{列ベクトル}は一回ごとに$3$倍され,第二の\bookterm{column-vector}{列ベクトル}は変わらない.
したがって$k$回更新した結果は
\[
  A^k\begin{pmatrix}1\\1\end{pmatrix}
    =3^k\begin{pmatrix}1\\1\end{pmatrix},
  \qquad
  A^k\begin{pmatrix}1\\-1\end{pmatrix}
    =\begin{pmatrix}1\\-1\end{pmatrix}
\]
とすぐに分かる.

一般に,零でない\bookterm{column-vector}{列ベクトル}$\boldsymbol{u}$が
\[
  A\boldsymbol{u}=\lambda\boldsymbol{u}
\]
を満たすとき,\,$\boldsymbol{u}$を$A$の\textbf{\bookterm{eigenvector}{固有ベクトル}},\,$\lambda$を対応する\textbf{\bookterm{eigenvalue}{固有値}}という.
\bookterm{eigenvalue}{固有値}は,そのベクトルに対して一回の更新が何倍になるかを表す数である.
零ベクトル$\boldsymbol{0}=\begin{pmatrix}0\\0\end{pmatrix}$は,どんな$\lambda$でもこの式を満たしてしまうため,\bookterm{eigenvector}{固有ベクトル}から除く.
負の\bookterm{eigenvalue}{固有値}では矢印の向きが反転し,\bookterm{eigenvalue}{固有値}$0$では零ベクトルになる場合もある.

\subsubsection{ベクトルを二つの部分に分ける}

\bookterm{eigenvector}{固有ベクトル}が見つかると,ほかのベクトルへの作用も計算しやすくなる.
たとえば
\[
  \begin{pmatrix}1\\0\end{pmatrix}
  =\frac12\begin{pmatrix}1\\1\end{pmatrix}
   +\frac12\begin{pmatrix}1\\-1\end{pmatrix}
\]
である.
右辺を実際に足せば,上が$1$,下が$0$になる.
\sectionref*{節：一次独立・基底・次元}では,独立な解の\bookterm{combination}{線形結合}で\bookterm{initial}{初期条件}を合わせた.
ここでは二つの独立な\bookterm{column-vector}{列ベクトル}を\bookterm{basis}{基底}に選び,元のベクトルを同じように\bookterm{combination}{線形結合}で表している.
\bookterm{sequence}{数列}全体を分けた見方を,今度は更新前の状態の分解に使っているのである.
二つの部分を別々に$k$回更新すれば
\[
  \begin{aligned}
    A^k\begin{pmatrix}1\\0\end{pmatrix}
      &=\frac{3^k}{2}\begin{pmatrix}1\\1\end{pmatrix}
        +\frac12\begin{pmatrix}1\\-1\end{pmatrix}\\
      &=\begin{pmatrix}\frac{3^k+1}{2}\\\frac{3^k-1}{2}\end{pmatrix}.
  \end{aligned}
\]
大きな累乗を一回ずつ計算する代わりに,定数倍になる二つの部分を追えばよい.
本編で等比型に帰着させると計算が簡単になったのと同じように,\bookterm{matrix}{行列}でも一定の倍率で変わる部分を見つけることが鍵になる.

一般に,二つの異なる\bookterm{eigenvalue}{固有値}$\lambda_1,\lambda_2$に対応する\bookterm{eigenvector}{固有ベクトル}$\boldsymbol{u}_1,\boldsymbol{u}_2$があれば,どんな二成分のベクトルも
\[
  \boldsymbol{v}=c\boldsymbol{u}_1+d\boldsymbol{u}_2
\]
と表せる.
異なる\bookterm{eigenvalue}{固有値}に対応する二つの\bookterm{eigenvector}{固有ベクトル}は,互いの定数倍にならないからである.
係数$c,d$は,上下の成分を比較した連立方程式で求められる.
その後は
\[
  A^k\boldsymbol{v}
  =c\lambda_1^k\boldsymbol{u}_1+d\lambda_2^k\boldsymbol{u}_2
\]
で計算できる.
$k=0$のときは元の分解式を読み,\,$k\ge1$で\bookterm{eigenvalue}{固有値}が$0$なら対応する部分は零になる.

\subsection{行ベクトルと左固有ベクトル}

数を横に並べた$L=(\alpha,\beta)$を\textbf{\bookterm{row-vector}{行ベクトル}}という.
\bookterm{column-vector}{列ベクトル}$\boldsymbol{\ell}$を横に並べ直す操作を\textbf{\bookterm{transpose}{転置}}といい,\,$\boldsymbol{\ell}^{\mathsf T}$と書く.
たとえば$\boldsymbol{\ell}=\begin{pmatrix}\alpha\\\beta\end{pmatrix}$なら,\,$\boldsymbol{\ell}^{\mathsf T}=(\alpha,\beta)$である.
\bookterm{row-vector}{行ベクトル}と\bookterm{column-vector}{列ベクトル}の積は,対応する成分を掛けて足す.
\[
  (\alpha,\beta)\begin{pmatrix}x\\y\end{pmatrix}
  =\alpha x+\beta y.
\]
\bookterm{row-vector}{行ベクトル}を\bookterm{matrix}{行列}に掛けるときは,\bookterm{matrix}{行列}の各列との積を横に並べる.
たとえば,先ほどの$A$について
\[
  \begin{aligned}
    (1,\,1)A&=(3,\,3)=3(1,\,1),\\
    (1,-1)A&=(1,-1)
  \end{aligned}
\]
である.
一般に,零でない\bookterm{row-vector}{行ベクトル}$L$が$LA=\gamma L$を満たすとき,\,$L$を\textbf{\bookterm{left-eigen}{左固有ベクトル}}という.
この場合,\bookterm{column-vector}{列ベクトル}$\boldsymbol{v}$から作った数$L\boldsymbol{v}$は,更新後には
\[
  L(A\boldsymbol{v})=(LA)\boldsymbol{v}
  =\gamma L\boldsymbol{v}
\]
となる.
たとえば$L=(1,\,1)$なら,\,$x+y$という和が更新によって$3$倍される.
\bookterm{eigenvector}{固有ベクトル}は定数倍になる\bookterm{column-vector}{列ベクトル}を表し,\bookterm{left-eigen}{左固有ベクトル}は定数倍になる量を取り出す係数を表す.

\subsection{固有値を求める条件：行列式}

\bookterm{eigenvector}{固有ベクトル}の成分を比較すれば,\bookterm{eigenvalue}{固有値}を求める条件が得られる.
その条件を一つの式にまとめるため,\bookterm{matrix}{行列}式という記号を導入しよう.
二行二列の場合には
\[
  \det\begin{pmatrix}a&b\\c&d\end{pmatrix}=ad-bc
\]
と定める.
\bookterm{matrix}{行列}そのものではなく,四つの成分から計算した一つの数である.

二つの式$ax+by=0,\ cx+dy=0$について,\,$ad-bc\ne0$なら解は$x=y=0$だけになる.
たとえば,第一の式を$d$倍したものから第二の式を$b$倍したものを引くと,\,$(ad-bc)x=0$となる.
同様に$(ad-bc)y=0$も得られる.
逆に$ad-bc=0$なら,零でない解が存在する.
実際,\,$(a,b)\ne(0,\,0)$なら$(x,y)=(b,-a)$が両方の式を満たす.
$(a,b)=(0,\,0)$の場合も,残る一つの式から零でない解を選べる.

\bookterm{eigenvalue}{固有値}の式$A\boldsymbol{u}=\lambda\boldsymbol{u}$は
\[
  (\lambda I-A)\boldsymbol{u}=\boldsymbol{0}
\]
と書き直せる.
\bookterm{matrix}{行列}の引き算は対応する成分ごとに行うので,\,$A=\begin{pmatrix}p&q\\r&s\end{pmatrix}$なら
\[
  \lambda I-A
  =\begin{pmatrix}\lambda-p&-q\\-r&\lambda-s\end{pmatrix}.
\]
零でない解が存在する条件は
\[
  \begin{aligned}
    \det(\lambda I-A)
      &=(\lambda-p)(\lambda-s)-qr\\
      &=\lambda^2-(p+s)\lambda+(ps-qr)=0.
  \end{aligned}
\]
これを\bookterm{matrix}{行列}$A$の\textbf{\bookterm{characteristic}{特性方程式}}という.
この方程式を解いて$\lambda$を求め,成分の連立方程式を解けば,対応する\bookterm{eigenvector}{固有ベクトル}も見つかる.
たとえば$A=\begin{pmatrix}2&1\\1&2\end{pmatrix}$では
\[
  \lambda^2-4\lambda+3
  =(\lambda-3)(\lambda-1)=0
\]
だから,\bookterm{eigenvalue}{固有値}は$3,\,1$である.
これは数値例で確かめた倍率と一致する.

\bookterm{recurrence}{漸化式}では,次の項を計算するために記録しておく値の組を\textbf{状態}と呼ぶ.
また,\bookterm{matrix}{行列}式が$0$でないときは,更新後の二つの値から更新前の二つの値を一意に求め直せる.
このように更新を逆向きに戻せる\bookterm{matrix}{行列}を\textbf{可逆な\bookterm{matrix}{行列}}という.

次節では,これらの計算を本編の\bookterm{recurrence}{漸化式}に当てはめていこう.
特に,本編で選んだ「何を引くか」「どのような比や和を作るか」が,\bookterm{matrix}{行列}では何を選ぶことに対応するのかを考える.
